\documentclass{oxmathproblems}

\printanswers

\course{Understanding Analysis, 2nd Edition: Stephen Abbott}
\sheettitle{Chapter 3 - Section 3.3 - Compact Sets} %can l7ave out if no title per sheet

\begin{document}
\begin{questions}

%------------------------- Problem 1 -------------------------
\miquestion Show that if $K$ is compact and nonempty, then sup $K$ and inf $K$ both exist and are elements of $K$.
\begin{solution}
  $K$ is a bounded subset of $\textbf{R}$, so sup $K$ and inf $K$ both exist by the Axiom of Completeness. They're both
  elements of $K$ because $K$ is closed (see exercise 3.2.4).
\end{solution}

%------------------------- Problem 2 -------------------------
\miquestion Decide which of the following sets are compact. For those that are not compact, show how Definition 3.3.1 breaks down. In other words, give an example
of a sequence contained in the given set that does not possess a subsequence converging to a limit in the set.
\begin{enumerate}[label=(\alph*)]
  \item \textbf{N}.
  \item $\textbf{Q} \cap [0, 1]$.
  \item The Cantor set.
  \item $\{1 + \frac{1}{2^2} +  \frac{1}{3^2} + ... + \frac{1}{n^2} : n \in \textbf{N}\}$.
  \item $\{1, \frac{1}{2}, \frac{2}{3}, \frac{3}{4}, \frac{4}{5}, ...\}$.
\end{enumerate}

\begin{solution}
  \begin{enumerate}[label=(\alph*)]
    \item Not compact because it's not bounded. Example sequence: $(1, 2, 3, 4, 5, ...)$. Any subsequence diverges to infinity.
    \item Not compact because it's not closed. Take an irrational $x \in [0, 1]$, say $x = \sqrt{2}/2$. There is a sequence of rationals bounded
    by [0, 1] that converges to $x$ (by density of the rationals)*. Any subsequence also converges to $x$, but $x \notin \textbf{Q}$.

    *or an explicit sequence by Heron's method (Exercise 2.4.5): let $x_{1} = \frac{1}{2}$ and $x_{n+1}=\frac{1}{2}(x_{n}+\frac{1}{2x_{n}})$.

    \item Compact. The Cantor set is clearly bounded by 0 and 1. By exercise 3.2.6(e), it's closed.

    \item Not compact because it's not closed (exercise 3.2.3 (d)). Example sequence: $(s_{n} = \sum_{i=1}^{n}\frac{1}{i^{2}})$. The sequence of partial sum
    converges to a limit $L$ (example 2.4.4), so any subsequence converges $L$. But $L$ is not one of the partial sums, so no subsequence
    converges to a point in the set.

    \item Compact. Instead of proving it's closed and bounded, let's use the (equivalent) Definition 3.3.1. Take any (not necessarily convergent) sequence $(x_{n})$
    from the set. We need to prove that $(x_{n})$ contains a subsequence that converges to a limit that's an element of the set.

    If $(x_{n})$ contains an element that appears infinitely many times then it's trivial. Pick that element as a subsequence.
    
    Otherwise, $(x_{n})$ contains infinitely many distinct values. These distinct values must include fractions $\frac{m}{m+1}$ with arbitrarily large $m$.
    Choose them with increasing $m$, giving a subsequence converging to 1, and 1 is in the set.
  \end{enumerate}
\end{solution}

%------------------------- Problem 3 -------------------------
\miquestion Prove the converse of Theorem 3.3.4 by showing that if a set $K \subseteq \textbf{R}$ is closed and bounded, then it is compact (every sequence has a convergent
subsequence with a limit in $K$).

\begin{solution}
  Take any sequence in $K$ (not necessarily convergent). Since $K$ is bounded, this sequence is also bounded. Every sequence has a monotone subsequence
  (Exercise 2.5.8) and this subsequence is also bounded. By Monotone Convergence, it converges and since $K$ is closed, the limit is in $K$.

  *later improvement: once we establish that the sequence is bounded, invoke Bolzano-Weierstrass directly to conclude the existence of a convergent subsequence.
\end{solution}

%------------------------- Problem 4 -------------------------
\miquestion Skip.

%------------------------- Problem 5 -------------------------
\miquestion Short proof if true, counterexample if false.
\begin{enumerate}[label=(\alph*)]
  \item The arbitrary intersection of compact sets is compact.
  \item The arbitrary union of compact sets is compact.
  \item Let $A$ be arbitrary and let $K$ be compact. Then the intersection $A \cap K$ is compact.
  \item If $F_{1} \supseteq F_{2} \supseteq F_{3} \supseteq ...$ is a nested sequence of nonempty closed sets, then the intersection $\bigcap_{n=1}^{\infty} F = \emptyset$.
\end{enumerate}

\begin{solution}
  \begin{enumerate}[label=(\alph*)]
    \item True. A set is compact if and only if it's closed and bounded. The intersection of an arbitrary collection of closed sets is closed (Theorem 3.2.14 (ii)). The intersection
    of an arbitrary collection of bounded sets is bounded.

    \item False. Cover the positive real line with $\{[0, 1] \cup [1, 2] \cup [2, 3] \cup ...\}$. A closed interval is compact, so this is a union of compact sets. But the resulting
    union is unbounded.

    What would be true is that \textit{finite} union of compact sets is compact.

    \item False. Let $A = (0, 1)$ and $K = [0, 1]$. The intersection $(0, 1)$ is not closed. (note: it's not enough to claim (0, 1) is open, as a set can be both open and closed).
    
    \item False. This is the same as Exercise 3.2.6(b), asking whether the Nested Interval Property is true for closed sets. 
    
    Counterexample: let $F_{n} = [n, \infty)$. The intersection is empty because
    any number $x$ is not in $F_{\lceil x \rceil + 1}$.

    As per Theorem 3.3.5, what is needed for this to be true is for each set to be \textit{compact}, not just closed.
  \end{enumerate}
\end{solution}

%------------------------- Problem 6 -------------------------
\miquestion This exercise is meant to illustrate the point made in the opening paragraph in Section 3.3. Verify that the following three statements are true if
every blank is filled in with the word "finite". Which are true if every blank is filled in with "compact". Which are true if filled in with "closed".

To quote the opening paragraph: "\textit{The central challenge in analysis is to exploit the power of the mathematical
infinite—via limits, series, derivatives, integrals, etc.—without falling victim to
erroneous logic or faulty intuition. A major tool for maintaining a rigorous
footing in this endeavor is the concept of compact sets. In ways that will become
clear, especially in our upcoming study of continuous functions, employing
compact sets in a proof often has the effect of bringing a finite quality to the
argument, thereby making it much more tractable.}"

\begin{enumerate}[label=(\alph*)]
  \item Every \underline{\phantom{blank}} set has a maximum.
  \item If $A$ and $B$ are \underline{\phantom{blank}}, then $A + B = \{a + b: a \in A, b \in B\}$ is also \underline{\phantom{blank}}.
  \item If $\{A_{n} : n \in \textbf{N}\}$ is a collection of \underline{\phantom{blank}} sets with the property that every finite subcollection has a nonempty intersection, then
  $\bigcap_{n=1}^{\infty} A_{n}$ is nonempty as well.
\end{enumerate}

\begin{solution}
  \begin{enumerate}[label=(\alph*)]
    \item That finite sets have a maximum is obvious.
    
    It's also true for compact sets. This is basically Exercise 3.3.1: the supremum of a compact set is an element of the set.
    If a supremum is part of a set, it must also be the maximum.

    It's \textit{not} necessarily true for closed sets, as closed sets can be unbounded e.g $[1, \infty)$.

    \item Obvious for finite sets.
    
    True for compact sets. I got help online for this and I'm annoyed. But anyway, here's what a plausible thought process could look like.

    It looks true intuitively, so let's try to prove it. If it doesn't work, the blocker may suggest how to prove it wrong.

    We need to prove that for any convergent sequence $(a_{n}+b_{n}) \in A+B$, the limit $L$ is in $A+B$. Now, we'll definitely need to invoke
    $A$ and $B$'s compactness i.e. produce a convergent sequence for each. But it can't just be any random sequence; it needs to come from (i.e.
    be a subsequence of) $(a_{n}+b_{n})$. But before we produce such a subsequence, let's verify whether it'd actually help. The answer is yes,
    because if we have a subsequence $(a_{n_{k}} + b_{n_{k}})$ such that $(a_{n_{k}})$ converges to $L_{a}$ and $(b_{n_{k}})$ to $L_{b}$, then by
    the Algebraic Limit Theorem $L_{a} + L_{b} = L$ and since $A$ and $B$ are closed, this implies $L_{a} \in A$, $L_{b} \in B$, hence $L = L_{a}+L_{b} \in A+B$,
    finishing the proof.  

    Now to produce the subsequence. It's easy to create a subsequence $(a_{n_{k}})$ that converges. Since $A$ is compact, any sequence in it has a
    convergent subsequence, so just take it out of $(a_{n}+b_{n})$. Likewise, we can create a separate subsequence $(b_{n_{j}})$ that converges. But
    the blocker here is that they have different indices! E.g. if the convergent subsequence for $A$ contains $a_{3}$ but not $a_{5}$, whereas the one for
    $B$ contains $b_{5}$ but not $b_{3}$, we can't create a \textit{single} subsequence that includes both $a_{3}+b_{3}$ and $a_{5}+b_{5}$, because then
    one might not converge.

    And here's the trick that I failed to come up with. Instead of creating a separate subsequence out of the single sequence $(a_{n}+b_{n})$, we can instead
    first take the subsequence $(a_{n_{k}}+b_{n_{k}})$ for $A$, then \textit{reuse this} to create the subsequence for $B$! So out of $(a_{n_{k}}+b_{n_{k}})$,
    pick a subsequence $(a_{n_{k_{j}}} + b_{n_{k_{j}}})$ such that $B$ converges.

    The remaining important question is: how to come up with this trick if we missed it the first time? I don't know the answer. I was too stubborn on trying
    to create both subsequences out of a single sequence. Maybe the lesson here is that whenever we're given a problem about satisfying
    multiple constraints (not necessarily convergence, but just constraints in general), instead of trying to satisfy all constraints simultaneously, try to 
    satisfy them by incrementally restricting the solution step by step. Are there other problems outside Real Analysis with this sort of trick?
    
    Or maybe this is one of those tricks that's easy to miss the first time, but will come up over and over again
    until it doesn't seem magic anymore.

    False for closed sets. Counterexample: $A = \{1 + \frac{1}{2}, 2 + \frac{1}{3}, 3 + \frac{1}{4}, ...\}$ and $B = \{-1, -2, -3, ...\}$. Then $A+B$ contains the
    sequence $(1/n)$ that converges to 0 but $0 \notin A+B$. 

    \item TODO.
  \end{enumerate}
\end{solution}

%------------------------- Problem 7 -------------------------
\miquestion As some more evidence of the surprising nature of the Cantor set, follow these steps to show that the sum $C + C = \{x+y: x, y \in C\}$
is equal to the closed interval $[0, 2]$. (Keep in mind that $C$ has zero length and contains no intervals).

Because $C \subseteq [0, 1]$, $C + C \subseteq [0, 2]$, so we only need to prove the reverse inclusion $[0, 2] \subseteq \{x+y: x, y \in C\}$. Thus,
given $s \in [0, 2]$, we must find two elements $x, y$ satisfying $x + y = s$.

\begin{enumerate}[label=(\alph*)]
  \item Show that there exist $x_{1}, y_{1} \in C_{1}$ for which $x_{1} + y_{1} = s$. Show in general that, for an arbitrary $n \in \textbf{N}$, we can
  always find $x_{n}, y_{n} \in C_{n}$ for which $x_{n}+y_{n}=s$.
  \item Keeping in mind that the sequences $(x_{n})$ and $(y_{n})$ do not necessarily converge, show how they can nevertheless be used to produce the desired
  $x$ and $y$ in $C$ satisfying $x+y=s$.
\end{enumerate}

\begin{solution}
  Let's first really appreciate why this is such a surprising result. As the book says, the Cantor set is really a "Cantor dust"; just a set of points,
  \textit{without any intervals}, zero length and is very sparse. How does summing two such sets result in a \textit{solid interval}, without gaps! Just as good as summing up
  the two solid intervals $[0, 1] + [0, 1]$. It's visually highly unintuitive.

  \begin{enumerate}[label=(\alph*)]
    \item There are much more elegant proofs at the internet. But here's my own ugly attempt.
    
    The way the problem is phrased hints at induction. The base case is easy: $C_{0} = [0, 1]$ and clearly $C_{0} + C_{0} = [0, 2]$. Now, let's see
    how to solve the problem for $C_{1}$ first, and then generalize later.

    Suppose we have $x, y$ both in $[0, 1]$ such that $x+y=s$. Can we reuse this $x, y$ for $C_{1}$? Well, the only problem is if one or both of $x, y$ lies
    in the newly created gap $(\frac{1}{3}, \frac{2}{3})$. But it should not be too difficult to intuitively see that we just need to shift $x$ and $y$ to be
    inside the two intervals in $C_{1}$ and this is always possible.
    
    More formally, wlog $x \leq y$:
    \begin{enumerate}
      \item If $x, y \in ([0, \frac{1}{3}] \cup [\frac{2}{3}, 1])$, then all good.
      \item If \textit{both} $x, y$ lie inside the gap $(\frac{1}{3}, \frac{2}{3})$, shift them to $x'=x-\frac{1}{3}$ and $y'=y+\frac{1}{3}$. Since the length
      of the gap is less than $\frac{1}{3}$, both $x', y'$ must now lie outside the gap and since $x, y$ being outside the gap implies $x > \frac{1}{3}$ and
      $y < \frac{2}{3}$, then $x' = x-\frac{1}{3} > \frac{1}{3}-\frac{1}{3} = 0$ and similarly $y' < 1$. We also have $x'+y' = s$. 
      
      But really, the visualization is much easier than the formal proof.

      \item If only one of $x, y$ lies inside the gap, wlog $y$ is in the gap (the other case of $x$ in the gap can be handled similarly). 
      
      First, shift to $x'= x-x = 0$ and $y'= y + x$. Since $y$ was in the gap and $x \leq \frac{1}{3}$, $y' < \frac{2}{3}+\frac{1}{3} = 1$, so $y'$ is still
      in $C_{1}$. If $y'$ is now inside the interval $[\frac{2}{3}, 1]$, then done.
     
      Otherwise, $y'$ is still in the gap. Since $x'=0$ and $y' < \frac{2}{3}$, their midpoint $m$ must be less than $\frac{1}{3}$. Shift both the $m$,
      and they're now both outside the gap.
    \end{enumerate}

    All this casework proves that $C_{1}+C_{1} = [0, 2]$, by reusing all the $x,y$ from $C_{0}$ and shifting them as needed. Now, the general case.

    Suppose $C_{n-1}+C_{n-1}=[0, 2]$. For the inductive case, we need to prove that $C_{n}+C_{n}=[0, 2]$. Given $s \in [0, 2]$, let $x, y \in C_{n-1}$ such
    that $x+y=s$. Note that for $n > 1$, the length of the gaps in $C_{n}$ are not all the same e.g. in $C_{2}$ there are gaps with length $\frac{1}{9}$
    and the middle gap with length $\frac{1}{3}$. In general, the smallest gap is the same length as the length of an interval in $C_{n}$ which is $\frac{1}{3^{n}}$.
    But since $x, y \in C_{n-1}$, if they lay inside a gap in $C_{n}$, it must be in the smallest gap of length $\frac{1}{3^{n}}$, because $C_{n-1}$ itself doesn't
    contain any element in the bigger gap (see the visualization in page 86, and note how $C_{2}$ doesn't contain any element in the big gap, so if any elements
    from $C_{2}$ ends up in a gap in $C_{3}$, it must be in the small gap).

    Anyway, just shift them like above, but this time shift with length $\frac{1}{3^{n}}$ It's just tedious technicality, I'm not in the mood to write it all out.
  
    \item TODO.
  \end{enumerate}

\end{solution}

\end{questions}
\end{document}
